% ===========================================================================
%  Capítulo: El lenguaje de dominio específico  — archivo autocontenido
% ===========================================================================
\documentclass[12pt,a4paper]{report}
\usepackage[utf8]{inputenc}
\usepackage[T1]{fontenc}
\usepackage[spanish]{babel}
\usepackage{geometry}
\geometry{left=3cm,right=2.5cm,top=3cm,bottom=2.5cm}
\usepackage{listings}
\usepackage{xcolor}
\usepackage{amsmath}
\usepackage{booktabs}
\usepackage{array}
\usepackage{parskip}
\usepackage{hyperref}
\usepackage{titlesec}
\usepackage{enumitem}

\definecolor{lispkw}{rgb}{0.0,0.35,0.65}
\definecolor{lispstr}{rgb}{0.13,0.55,0.13}
\definecolor{lispcomment}{rgb}{0.5,0.5,0.5}
\definecolor{lispbg}{rgb}{0.97,0.97,0.99}
\definecolor{pybg}{rgb}{0.97,0.99,0.97}
\definecolor{nodeborder}{rgb}{0.3,0.45,0.7}
\definecolor{sectioncolor}{rgb}{0.15,0.25,0.5}

\lstdefinelanguage{CommonLisp}{
  keywords={defclass, defmacro, defun, defparameter, defvar, progn, let, let*,
            loop, for, collect, append, list, when, unless, if, cond,
            load, format, lambda, provide, values,
            def-table, conditional-rendering, exists, col, trange,
            countif, counta, sum-range, str, nav, ultima-fila, primera-fila,
            libro, hoja, tabla, xl-generate, xl-run-generated},
  keywordstyle=\color{lispkw}\bfseries,
  sensitive=false,
  comment=[l]{;;},
  commentstyle=\color{lispcomment}\itshape,
  morestring=[b]",
  stringstyle=\color{lispstr},
  basicstyle=\ttfamily\small,
  backgroundcolor=\color{lispbg},
  frame=single, framerule=0.4pt,
  rulecolor=\color{nodeborder!60},
  breaklines=true, showstringspaces=false,
  tabsize=2, xleftmargin=1em, xrightmargin=0.5em,
  aboveskip=0.8em, belowskip=0.5em,
}
\lstdefinestyle{pythonstyle}{
  language=Python,
  basicstyle=\ttfamily\small,
  backgroundcolor=\color{pybg},
  keywordstyle=\color{lispkw}\bfseries,
  stringstyle=\color{lispstr},
  commentstyle=\color{lispcomment}\itshape,
  frame=single, framerule=0.4pt,
  rulecolor=\color{nodeborder!60},
  breaklines=true, showstringspaces=false,
  tabsize=4, xleftmargin=1em, xrightmargin=0.5em,
  aboveskip=0.8em, belowskip=0.5em,
}
\lstset{language=CommonLisp}

\hypersetup{
  pdftitle={El Lenguaje de Dominio Espec\'{\i}fico},
  pdfauthor={Jose Morales Lazo},
  colorlinks=true,
  linkcolor=sectioncolor,
}

\titleformat{\chapter}[display]
  {\normalfont\Large\bfseries\color{sectioncolor}}
  {\chaptertitlename\ \thechapter}{16pt}{\Huge}
\titleformat{\section}
  {\normalfont\large\bfseries\color{sectioncolor}}{\thesection}{1em}{}
\titleformat{\subsection}
  {\normalfont\normalsize\bfseries\color{sectioncolor!80}}{\thesubsection}{1em}{}

\begin{document}

\chapter{El Lenguaje de Dominio Espec\'{\i}fico}
\label{cap:dsl}

% ---------------------------------------------------------------------------
\section{Visi\'on general}
\label{sec:vision-general}
% ---------------------------------------------------------------------------

Numerosos problemas cotidianos producen informaci\'on que se organiza
naturalmente en forma de tabla: filas que representan entidades y columnas
que describen sus atributos, relaciones o restricciones.  Los horarios son
un ejemplo paradigm\'atico de esta clase de datos---franjas horarias,
participantes, recursos y conflictos pueden expresarse como tablas
relacionadas---pero la misma estructura aparece en inventarios, resultados
acad\'emicos, planificaci\'on de proyectos y cualquier dominio donde la
informaci\'on admita una organizaci\'on tabular.

El objetivo del lenguaje de dominio espec\'{\i}fico (DSL) desarrollado en
este trabajo es ofrecer un recurso de descripci\'on para esa clase de datos:
un vocabulario declarativo con el que el programador expresa \emph{qu\'e}
contiene cada tabla---columnas, f\'ormulas derivadas, relaciones entre
tablas, reglas de presentaci\'on---sin ocuparse de \emph{c\'omo} se
materializa esa descripci\'on en un formato concreto.

El DSL est\'a implementado como un conjunto de funciones y macros de
Common~Lisp que construyen, en tiempo de carga, un \emph{\'arbol de sintaxis
abstracta} (AST) cuyos nodos son instancias de clases CLOS.  La canalizaci\'on
de compilaci\'on comprende tres etapas claramente separadas:

\begin{enumerate}
  \item \textbf{Expansi\'on de macros.}  Cada forma del DSL
    (p.\,ej.\ \texttt{def-table}, \texttt{exists})se expande en una llamada al constructor
    CLOS del nodo AST correspondiente.  El resultado es un \'arbol cuyos
    nodos representan las entidades a partir de las cuales el sistema es
    capaz de generar las representaciones visuales de los datos.

  \item \textbf{Generaci\'on de c\'odigo.}  El \'arbol se recorre por el
    backend, que traduce cada nodo a las instrucciones propias del formato de
    salida: expresiones de celda, reglas de presentaci\'on y directivas de
    construcci\'on del documento.  El backend es el \'unico punto que conoce
    los nombres de columna reales, los \'{\i}ndices de fila y los detalles
    de la API de destino.

  \item \textbf{Ejecuci\'on.}  El programa generado se ejecuta, produciendo
    el archivo de salida con los datos, las f\'ormulas derivadas y las reglas
    de presentaci\'on incrustadas.
\end{enumerate}

Un principio central de dise\~no es la \emph{separaci\'on entre intenci\'on
e implementaci\'on}: el DSL expresa la sem\'antica del problema
(``resaltar los profesores que tienen un conflicto de horario'') y el backend
decide el mecanismo concreto seg\'un el formato de destino elegido.  El
programador del DSL no escribe ninguna instrucci\'on espec\'{\i}fica del
formato de salida.

% ---------------------------------------------------------------------------
\section{\'Arbol de sintaxis abstracta}
\label{sec:ast}
% ---------------------------------------------------------------------------

El \'arbol de sintaxis abstracta (AST) es la representaci\'on intermedia que
separa la capa del DSL de la capa del generador de c\'odigo.  Cada concepto
del problema---una tabla, una expresi\'on de celda, un cuantificador
existencial---queda capturado como un nodo AST con atributos propios.  El
backend opera \'unicamente sobre este \'arbol; no interpreta texto ni
inspecciona las formas del DSL.

Todos los nodos del AST se definen con la macro \texttt{defclass*}, que
genera simult\'aneamente una clase CLOS \texttt{clase-xl-\textit{nombre}} y
una funci\'on constructora \texttt{xl-\textit{nombre}} con par\'ametros por
clave.  Por ejemplo, el nodo del cuantificador existencial se define y
construye de la siguiente manera:

\begin{lstlisting}[language=CommonLisp,
  caption={Patron de definicion de nodo AST con defclass*.}]
;; defclass* genera la clase y el constructor xl-expr-exists:
(defclass* xl-expr-exists ()
           (bind-var domain match-keys overlap-cols self-in-cols))

;; El constructor acepta los slots como :keywords:
(xl-expr-exists :bind-var     'r
                :domain       (xl-domain-rows :table-id nil)
                :match-keys   '(dia hora)
                :overlap-cols '(tutor oponente presidente vocal secretario)
                :self-in-cols  nil)
\end{lstlisting}

Los nodos del AST se agrupan en seis categor\'{\i}as seg\'un su papel en
el \'arbol: estructura del workbook, expresiones de celda, navegaci\'on de
fila, cuantificaci\'on existencial, expresiones cross-sheet din\'amicas y
formato condicional con posicionamiento relativo.

% ............................................................................
\subsection{Estructura del workbook}
\label{sec:ast-estructura}
% ............................................................................

Los nodos de estructura representan los contenedores del libro de c\'alculo.
Son los nodos de m\'as alto nivel del \'arbol: el workbook agrupa hojas, las
hojas agrupan regiones y las regiones agrupan tablas.  Su jerarqu\'{\i}a
refleja directamente la estructura visible del archivo \texttt{.xlsx}, de
modo que el programador puede razonar sobre la disposici\'on del libro
usando los mismos t\'erminos que usar\'{\i}a al abrirlo en Excel.

\begin{table}[h!]
\centering
\caption{Nodos de estructura del workbook.}
\label{tab:ast-estructura}
\renewcommand{\arraystretch}{1.3}
\begin{tabular}{p{3.6cm} p{4.8cm} p{5.4cm}}
\hline
\textbf{Nodo} & \textbf{Atributos} & \textbf{Papel} \\
\hline
\texttt{xl-workbook}
  & \texttt{name}, \texttt{sheets}
  & Ra\'{\i}z del \'arbol.  \texttt{name} es el nombre del
    archivo \texttt{.xlsx}; \texttt{sheets} es la lista de
    hojas. \\
\hline
\texttt{xl-sheet}
  & \texttt{name}, \texttt{regions},\newline
    \texttt{fixed-expressions}
  & Una hoja del libro.  Contiene regiones horizontales y
    expresiones en posici\'on relativa fuera de las tablas. \\
\hline
\texttt{xl-region}
  & \texttt{tables}
  & Agrupa tablas colocadas una junto a otra en la misma
    hoja, separadas por una columna en blanco. \\
\hline
\texttt{xl-table}
  & \texttt{id}, \texttt{col-names},\newline
    \texttt{headers}, \texttt{contenido},\newline
    \texttt{computed}, \texttt{style-rules},\newline
    \texttt{cell-height}, \texttt{cell-width}
  & Plantilla de tabla.  \texttt{col-names} enumera los
    s\'{\i}mbolos DSL de columna; \texttt{computed} asocia
    columnas a expresiones de f\'ormula; \texttt{style-rules}
    almacena las reglas de coloreado condicional. \\
\hline
\end{tabular}
\end{table}

% ............................................................................
\subsection{Expresiones de celda}
\label{sec:ast-expresiones}
% ............................................................................

Las expresiones de celda son los nodos que el backend traduce a cadenas de
f\'ormula Excel.  Abarcan tanto referencias a columnas de la misma tabla o
de otra (\texttt{xl-expr-column-ref}, \texttt{xl-table-range}) como
operaciones de hoja de c\'alculo concretas (\texttt{COUNTIF},
\texttt{COUNTA}, \texttt{SUM}) y literales de texto.  Gracias a esta
representaci\'on, el programador escribe \texttt{(countif rango criterio)}
en lugar de componer manualmente la cadena
\texttt{"COUNTIF(\$B\$4:\$F\$6,K4)"}.

\begin{table}[h!]
\centering
\caption{Nodos de expresion de celda.}
\label{tab:ast-expresiones}
\renewcommand{\arraystretch}{1.3}
\begin{tabular}{p{4.2cm} p{3.8cm} p{5.8cm}}
\hline
\textbf{Nodo} & \textbf{Atributos} & \textbf{Papel} \\
\hline
\texttt{xl-expr-column-ref}
  & \texttt{name}, \texttt{context}
  & Referencia a la columna \texttt{name} en la fila actual.
    Produce la letra Excel de esa columna m\'as el n\'umero
    de fila relativo. \\
\hline
\texttt{xl-range}
  & \texttt{from-col}, \texttt{to-col}
  & Rango de columnas contiguas de la tabla actual. \\
\hline
\texttt{xl-table-range}
  & \texttt{table-id},\newline
    \texttt{from-col}, \texttt{to-col}
  & Rango en una tabla espec\'{\i}fica de la regi\'on
    (referencia cross-table).  Lo produce la forma
    \texttt{trange}. \\
\hline
\texttt{xl-expr-countif}
  & \texttt{count-range},\newline\texttt{criteria}
  & Formula \texttt{COUNTIF} sobre un rango con criterio. \\
\hline
\texttt{xl-expr-counta}
  & \texttt{count-range}
  & Formula \texttt{COUNTA}: cuenta celdas no vac\'{\i}as. \\
\hline
\texttt{xl-expr-sum}
  & \texttt{count-range}
  & Formula \texttt{SUM} sobre un rango. \\
\hline
\texttt{xl-expr-string}
  & \texttt{value}
  & Literal de texto; produce una cadena fija en la celda. \\
\hline
\texttt{xl-expr-if}
  & \texttt{test}, \texttt{then}, \texttt{else}
  & Condicional \texttt{IF} de tres ramas. \\
\hline
\texttt{xl-expr-equals}
  & \texttt{a}, \texttt{b}
  & Igualdad entre dos expresiones. \\
\hline
\texttt{xl-expr-and}, \texttt{xl-expr-or}
  & \texttt{a}, \texttt{b}
  & Conjunci\'on y disyunci\'on l\'ogica. \\
\hline
\texttt{xl-expr-concat}
  & \texttt{a}, \texttt{b}
  & Concatenaci\'on de texto (\texttt{\&}). \\
\hline
\texttt{xl-expr-table-col-ref}
  & \texttt{table-id}, \texttt{name},\newline\texttt{context}
  & Referencia a una columna de una tabla espec\'{\i}fica de la regi\'on.
    Evita la ambig\"uedad cuando dos tablas comparten el mismo nombre de
    columna. \\
\hline
\texttt{xl-expr-param-ref}
  & \texttt{name}
  & Referencia a un par\'ametro de instancia almacenado en celda auxiliar. \\
\hline
\texttt{xl-expr-different}
  & \texttt{a}, \texttt{b}
  & Desigualdad: produce \texttt{a <> b} en la f\'ormula Excel. \\
\hline
\texttt{xl-expr-gt}, \texttt{xl-expr-lt}
  & \texttt{a}, \texttt{b}
  & Comparaci\'on estricta: mayor que y menor que. \\
\hline
\texttt{xl-expr-gte}, \texttt{xl-expr-lte}
  & \texttt{a}, \texttt{b}
  & Comparaci\'on no estricta: mayor o igual y menor o igual. \\
\hline
\texttt{xl-expr-add}, \texttt{xl-expr-subtract}
  & \texttt{a}, \texttt{b}
  & Suma y resta aritm\'etica. \\
\hline
\texttt{xl-expr-multiply}, \texttt{xl-expr-divide}
   & \texttt{a}, \texttt{b}
   & Multiplicaci\'on y divisi\'on aritm\'etica. \\
\hline
\texttt{xl-expr-promedio}
   & \texttt{cols}
   & Promedio de los valores de las columnas indicadas:
     suma todos los valores y divide entre el n\'umero de columnas. \\
\hline
\texttt{xl-expr-time-add}
  & \texttt{a}, \texttt{b}
  & Suma de tiempos como texto:
    \texttt{TEXT(TIMEVALUE(a)+(b/1440),"hh:mm")}. \\
\hline
\texttt{xl-expr-lookup}
  & \texttt{value-field},\newline\texttt{key-expr}
  & \texttt{VLOOKUP} en la tabla actual: busca \texttt{key-expr} y devuelve
    la columna \texttt{value-field}. \\
\hline
\end{tabular}
\end{table}

% ............................................................................
\subsection{Navegaci\'on de fila y control de visibilidad}
\label{sec:ast-navegacion-fila}
% ............................................................................

Ciertos c\'alculos necesitan referenciar la fila anterior o siguiente de la
misma tabla---por ejemplo, acumular un valor incremental o mostrar una
transici\'on---o bien producir una celda visualmente vac\'{\i}a bajo alguna
condici\'on.  Estos nodos encapsulan esa l\'ogica sin exponer \'{\i}ndices de
fila al programador del DSL.

\begin{table}[h!]
\centering
\caption{Nodos de navegaci\'on de fila y control de visibilidad.}
\label{tab:ast-navegacion}
\renewcommand{\arraystretch}{1.3}
\begin{tabular}{p{3.8cm} p{3.8cm} p{6.2cm}}
\hline
\textbf{Nodo} & \textbf{Atributos} & \textbf{Papel} \\
\hline
\texttt{xl-expr-previous-row}
  & \texttt{expr}
  & Eval\'ua \texttt{expr} referenciando la columna correspondiente en la
    fila anterior de la tabla. \\
\hline
\texttt{xl-expr-next-row}
  & \texttt{expr}
  & Eval\'ua \texttt{expr} referenciando la columna correspondiente en la
    fila siguiente de la tabla. \\
\hline
\texttt{xl-expr-first-row}
  & ---
  & Produce \texttt{TRUE} cuando la fila en evaluaci\'on es la primera fila
    de datos (\texttt{ROW()=first-row}). \\
\hline
\texttt{xl-expr-non-empty}
  & \texttt{expr}
  & Produce \texttt{TRUE} cuando \texttt{expr} no es vac\'{\i}a ni cero. \\
\hline
\texttt{xl-expr-show-nothing}
   & ---
   & Produce \texttt{""}: la celda queda visualmente vac\'{\i}a. \\
\hline
\texttt{xl-expr-empty}
   & ---
   & Vac\'{\i}o sem\'antico.  El backend decide la representaci\'on:
     \texttt{""} en Excel, \texttt{null} en otros formatos. \\
\hline
\end{tabular}
\end{table}

% ...............................................
\subsection{Cuantificaci\'on existencial}
\label{sec:ast-cuantificacion}
% ............................................................................

La cuantificaci\'on existencial es el mecanismo del DSL para expresar
condiciones de la forma ``existe alguna fila que cumpla ciertos criterios''.
Mantenerla como un nodo AST propio, y no como una expresi\'on booleana
expl\'{\i}cita, permite al backend generar la comprobaci\'on m\'as adecuada
seg\'un el contexto: una formula \texttt{SUMPRODUCT} vectorial en el caso de
Excel, u otra implementaci\'on para un backend diferente.

\begin{table}[h!]
\centering
\caption{Nodos de cuantificacion existencial.}
\label{tab:ast-cuantificacion}
\renewcommand{\arraystretch}{1.3}
\begin{tabular}{p{3.6cm} p{4.8cm} p{5.4cm}}
\hline
\textbf{Nodo} & \textbf{Atributos} & \textbf{Papel} \\
\hline
\texttt{xl-expr-exists}
  & \texttt{bind-var}, \texttt{domain},\newline
    \texttt{match-keys},\newline
    \texttt{overlap-cols},\newline
    \texttt{self-in-cols}
  & Cuantificador existencial estructurado.  Declara que debe
    existir una fila del dominio que cumpla las condiciones
    expresadas por los slots sem\'anticos.  El backend decide
    el mecanismo de comprobaci\'on seg\'un el tipo de dominio. \\
\hline
\texttt{xl-domain-rows}
  & \texttt{table-id}
  & Dominio de filas.  \texttt{table-id = nil} indica la
    tabla actual; un s\'{\i}mbolo identifica una tabla
    espec\'{\i}fica de la regi\'on (dominio cross-table). \\
\hline
\end{tabular}
\end{table}

% ............................................................................
\subsection{Expresiones cross-sheet din\'amicas}
\label{sec:ast-cross-sheet}
% ............................................................................

Cuando un libro contiene m\'ultiples hojas con estructura homog\'enea---por
ejemplo, una hoja por grupo o aula---es necesario agregar valores a trav\'es
de ellas.  A diferencia de \texttt{xl-expr-cross-sheet-ref}, que referencia
una hoja con nombre fijo, los nodos de esta secci\'on trabajan con una
\emph{variable de hoja} resuelta en tiempo de compilaci\'on mediante el
entorno din\'amico \texttt{*sheet-env*}, declarado por \texttt{collect-over}.

\begin{table}[h!]
\centering
\caption{Nodos de expresi\'on cross-sheet din\'amica.}
\label{tab:ast-cross-sheet}
\renewcommand{\arraystretch}{1.3}
\begin{tabular}{p{3.8cm} p{4.4cm} p{5.6cm}}
\hline
\textbf{Nodo} & \textbf{Atributos} & \textbf{Papel} \\
\hline
\texttt{xl-expr-cross-sheet-ref}
  & \texttt{sheet},\newline\texttt{cell-template}
  & Referencia a una celda en otra hoja con nombre fijo.
    \texttt{cell-template} puede incluir \texttt{\$\{row-num\}} para fila
    din\'amica. \\
\hline
\texttt{xl-expr-cross-cell}
  & \texttt{sheet}, \texttt{xcol},\newline\texttt{row}
  & Referencia a una celda en la hoja ligada a \texttt{sheet} en
    \texttt{*sheet-env*}.  \texttt{xcol} se resuelve a letra Excel;
    \texttt{row} puede ser un entero fijo o una expresi\'on. \\
\hline
\texttt{xl-expr-source-row}
  & \texttt{table-id},\newline\texttt{offset}
  & Fila de una tabla origen mapeada desde la fila l\'ogica actual.
    \texttt{offset} selecciona la subcelda dentro de una fila compuesta
    (0 = primera, 1 = segunda, \ldots). \\
\hline
\texttt{xl-expr-sheet-id}
  & \texttt{sheet}
  & Identificador can\'onico de la hoja ligada a \texttt{sheet} en
    \texttt{*sheet-env*}.  Cada backend decide c\'omo recuperarlo. \\
\hline
\texttt{xl-expr-collect-over}
  & \texttt{groups},\newline\texttt{sheet-var},\newline\texttt{body}
  & Itera \texttt{body} sobre cada hoja de \texttt{groups}, ligando
    \texttt{sheet-var}, y concatena los resultados como lista separada por
    comas mediante \texttt{SUBSTITUTE(TRIM(\ldots)," ",",")}. \\
\hline
\end{tabular}
\end{table}

% ............................................................................
\subsection{Formato condicional y posicionamiento}
\label{sec:ast-formato}
% ............................................................................

Estos nodos permiten expresar dos capacidades que van m\'as all\'a de los
datos de la tabla: el coloreado condicional de celdas y la colocaci\'on de
expresiones en posiciones relativas a las tablas.  La separaci\'on en nodos
propios garantiza que el DSL no exponga coordenadas absolutas: el backend
resuelve las posiciones finales en funci\'on del tama\~no real de los datos.

\begin{table}[h!]
\centering
\caption{Nodos de formato condicional y posicionamiento relativo.}
\label{tab:ast-formato}
\renewcommand{\arraystretch}{1.3}
\begin{tabular}{p{3.6cm} p{4.8cm} p{5.4cm}}
\hline
\textbf{Nodo} & \textbf{Atributos} & \textbf{Papel} \\
\hline
\texttt{xl-style-rule}
  & \texttt{rule-condition},\newline
    \texttt{target-columns}
  & Regla de coloreado condicional.  \texttt{rule-condition}
    es un nodo de expresi\'on booleana; \texttt{target-columns}
    enumera las columnas que se colorean cuando la condici\'on
    es verdadera. \\
\hline
\texttt{xl-anchor}
  & \texttt{table-id},\newline
    \texttt{col-name},\newline
    \texttt{anchor-type}
  & Punto de referencia dentro de una tabla.
    \texttt{anchor-type} es \texttt{:ultima-fila} o
    \texttt{:primera-fila}. \\
\hline
\texttt{xl-nav}
  & \texttt{anchor}, \texttt{steps}
  & Posici\'on calculada como un ancla m\'as una secuencia de
    pasos direccionales (arriba, abajo, izquierda, derecha). \\
\hline
\texttt{xl-sheet-fixed-expr}
  & \texttt{pos}, \texttt{expr}
  & Expresi\'on colocada en una posici\'on relativa de la hoja.
    \texttt{pos} es un nodo \texttt{xl-nav}; \texttt{expr} es
    el nodo de expresi\'on. \\
\hline
\end{tabular}
\end{table}

% ---------------------------------------------------------------------------
\section{Definici\'on de tablas: \texttt{def-table}}
\label{sec:def-table}
% ---------------------------------------------------------------------------

La forma central del DSL para definir datos tabulares es \texttt{def-table}.
Produce una \emph{plantilla} de tabla reutilizable: una descripci\'on
abstracta de las columnas, sus cabeceras, las f\'ormulas que se calculan a
partir de ellas y la regla de coloreado que se aplica sobre sus filas.  La
plantilla no contiene datos; se instancia con datos concretos al ensamblar
el workbook (secci\'on~\ref{sec:workbook}).

La sintaxis completa de \texttt{def-table} es:

\begin{lstlisting}[language=CommonLisp,
  caption={Sintaxis completa de def-table.}]
(def-table nombre-tabla
  ((col1 "Cabecera1") (col2 "Cabecera2") ...)  ; columnas
  :cell-height N       ; filas por fila logica (default 1)
  :cell-width  N       ; columnas por columna logica (default 1)
  :computed            ; formulas derivadas (opcional)
    ((colN (expresion-dsl)))
  :render              ; regla de coloreado (opcional)
    (conditional-rendering
      :condition (expresion-booleana)
      :target-columns (col1 col2 ...)))
\end{lstlisting}

Cada par \texttt{(col "Cabecera")} declara un nombre de columna DSL
(s\'{\i}mbolo) y el texto de su encabezado visual en la hoja de c\'alculo.
El orden de las columnas es el orden f\'{\i}sico en Excel; este orden es
relevante cuando se usan rangos contiguos con \texttt{trange}.

La clave \texttt{:computed} asocia columnas a expresiones DSL que el
compilador traduce a f\'ormulas Excel.  Los valores iniciales de esas
columnas en los datos de entrada se ignoran: la celda siempre mostrar\'a el
resultado de la f\'ormula calculada.

La clave \texttt{:render} acepta una forma \texttt{conditional-rendering}
que especifica bajo qu\'e condici\'on se colorean determinadas columnas.
El mecanismo concreto lo decide el backend seg\'un el tipo de condici\'on
declarada (secci\'on~\ref{sec:formato-condicional}).

El listado siguiente define la tabla de profesores del caso de uso
principal, incluyendo una columna calculada y una regla de coloreado
condicional:

\begin{lstlisting}[language=CommonLisp,
  caption={Tabla de profesores con columna computed y regla de coloreado.}]
(def-table profesores-table
  ((nombre "") (grado "") (defensas ""))
  :cell-height 1
  :cell-width  1
  :computed
    ;; Cuenta cuantas veces aparece este profesor en los cinco roles
    ((defensas (countif (trange defensas-table tutor secretario)
                        (col nombre))))
  :render
    (conditional-rendering
      :condition
        (exists (r :rows-of defensas-table)
          :self-in  (tutor oponente presidente vocal secretario)
          :matching (dia hora))
      :target-columns (nombre)))
\end{lstlisting}

La variable \texttt{r} en la forma \texttt{exists} es un alias documental
de la fila iterada, an\'alogo a un alias de tabla en SQL: facilita la
lectura del c\'odigo sin que el generador lo eval\'ue directamente.  La
regla de coloreado declara que la celda \texttt{nombre} debe resaltarse
cuando el profesor participa en alguna defensa con conflicto de horario; el
mecanismo lo decide el backend (secci\'on~\ref{sec:formato-condicional}).

% ---------------------------------------------------------------------------
\section{Expresiones de celda}
\label{sec:expresiones}
% ---------------------------------------------------------------------------

El DSL ofrece un conjunto de formas simb\'olicas para construir expresiones
de celda sin escribir cadenas de f\'ormula Excel directamente.  El
compilador resuelve los nombres de columna a sus letras reales, los rangos
a sus coordenadas absolutas y las operaciones a su sintaxis Excel; el
programador solo declara la intenci\'on.

La tabla~\ref{tab:formas-expresion} recoge las formas disponibles junto con
su traducci\'on.

\begin{table}[h!]
\centering
\caption{Formas de expresion del DSL y su traduccion Excel.}
\label{tab:formas-expresion}
\renewcommand{\arraystretch}{1.3}
\begin{tabular}{p{5.5cm} p{8.3cm}}
\hline
\textbf{Forma DSL} & \textbf{Significado / traduccion Excel} \\
\hline
\texttt{(col nombre)}
  & Valor de la columna \texttt{nombre} en la fila actual.
    Se resuelve a la letra Excel de esa columna m\'as el
    n\'umero de fila relativo, p.\,ej.\ \texttt{B4}. \\
\hline
\texttt{(trange tabla-id desde hasta)}
  & Rango que abarca desde la columna \texttt{desde} hasta la
    columna \texttt{hasta} de \texttt{tabla-id}, para todas
    las filas de datos.  P.\,ej.\ \texttt{\$B\$4:\$F\$6}. \\
\hline
\texttt{(countif rango criterio)}
  & Formula \texttt{COUNTIF(rango, criterio)}. \\
\hline
\texttt{(counta rango)}
  & Formula \texttt{COUNTA(rango)}: n\'umero de celdas
    no vac\'{\i}as. \\
\hline
\texttt{(sum-range rango)}
  & Formula \texttt{SUM(rango)}. \\
\hline
\texttt{(str "texto")}
  & Literal de texto; coloca la cadena directamente en la
    celda. \\
\hline
\texttt{(\_equals a b)}
  & Igualdad: produce \texttt{(a = b)} en la f\'ormula Excel. \\
\hline
\texttt{(\_and a b)}, \texttt{(\_or a b)}
  & Conjunci\'on y disyunci\'on:
    \texttt{AND(a,b)}, \texttt{OR(a,b)}. \\
\hline
\texttt{(nav ancla dir n \ldots)}
  & Posici\'on relativa: ancla m\'as pasos direccionales.
    Produce una referencia de celda absoluta. \\
\hline
\end{tabular}
\end{table}

La forma \texttt{trange} merece atenci\'on especial porque produce
\emph{referencias cross-table}: permite que una expresi\'on en
\texttt{profesores-table} cite columnas de \texttt{defensas-table} sin que
el programador conozca las letras de columna reales.  El compilador las
resuelve en funci\'on del orden de declaraci\'on de las tablas en la
regi\'on.

\begin{lstlisting}[language=CommonLisp,
  caption={Expresiones DSL y sus formulas Excel generadas.}]
;; (col nombre) en profesores-table se resuelve a la columna K, fila relativa
(col nombre)          ;; -> K4  (en la fila 4)

;; (trange defensas-table tutor secretario) -> rango absoluto de los cinco roles
(trange defensas-table tutor secretario)   ;; -> $B$4:$F$6

;; La formula COUNTIF resultante
(countif (trange defensas-table tutor secretario) (col nombre))
;; -> COUNTIF($B$4:$F$6, K4)
\end{lstlisting}

% ---------------------------------------------------------------------------
\section{El cuantificador existencial \texttt{exists}}
\label{sec:exists}
% ---------------------------------------------------------------------------

El cuantificador \texttt{exists} permite declarar condiciones de
\emph{existencia} sobre filas de una tabla.  La motivaci\'on de esta forma
es que preguntas del tipo ``existe otra defensa con el mismo horario'' o
``aparece este profesor en alguna defensa con conflicto'' son muy frecuentes
en problemas de planificaci\'on, pero expresarlas mediante disyunciones y
bucles expl\'{\i}citos produce c\'odigo verboso y dif\'{\i}cil de leer.

Con \texttt{exists} el programador declara la intenci\'on sem\'antica a alto
nivel; el backend se encarga de generar la comprobaci\'on correcta para cada
caso.

El cuantificador existe en dos variantes seg\'un el dominio de b\'usqueda.

% ............................................................................
\subsection{Variante same-table: \texttt{:all-rows}}
\label{sec:exists-all-rows}
% ............................................................................

Esta variante busca otra fila dentro de la \emph{misma} tabla que satisfaga
condiciones de coincidencia de valores y solapamiento en columnas de rol.
Se usa cuando se quiere detectar si una fila del propio conjunto tiene
alg\'un conflicto con otra fila del mismo conjunto; por ejemplo, dos
defensas programadas en el mismo d\'{\i}a y hora con al menos un integrante
de tribunal en com\'un.

\begin{lstlisting}[language=CommonLisp,
  caption={Cuantificador exists sobre la tabla actual.}]
(exists (r :all-rows)
  :matching    (dia hora)              ; mismos valores en ambas filas
  :any-overlap (tutor oponente presidente vocal secretario))
                                       ; algun valor compartido en estos slots
\end{lstlisting}

La variable \texttt{r} es el alias de la fila del dominio que se est\'a
comprobando.  La sem\'antica de la forma anterior es: ``existe otra fila de
esta tabla (distinta de la actual) cuya columna \texttt{dia} y cuya columna
\texttt{hora} coinciden con los de la fila actual, y que adem\'as comparte
al menos un valor en las columnas de rol''.

% ............................................................................
\subsection{Variante cross-table: \texttt{:rows-of}}
\label{sec:exists-rows-of}
% ............................................................................

Esta variante busca en las filas de \emph{otra} tabla, relacionando el
valor de la celda actual con columnas de esa tabla.  Se usa cuando el
elemento que se eval\'ua (p.\,ej.\ un profesor) y los datos de planificaci\'on
(p.\,ej.\ las defensas) residen en tablas distintas.

\begin{lstlisting}[language=CommonLisp,
  caption={Cuantificador exists sobre otra tabla (cross-table).}]
(exists (r :rows-of defensas-table)
  :self-in  (tutor oponente presidente vocal secretario)
               ; el valor de la celda actual aparece en alguno de estos slots
  :matching (dia hora))
               ; ese slot tiene mas de una defensa programada
\end{lstlisting}

La sem\'antica de la forma anterior es: ``existe alguna fila de
\texttt{defensas-table} en la que el valor de la celda actual (el nombre
del profesor) aparece en una de las cinco columnas de rol, y cuyo par
(\texttt{dia}, \texttt{hora}) tiene m\'as de una defensa programada''.
El nombre \texttt{r} identifica, en la lectura del c\'odigo, la fila de
\texttt{defensas-table} que se est\'a comprobando.

\medskip
\noindent
En ambas variantes el alias de fila es puramente documental: el backend
genera directamente una formula vectorial SUMPRODUCT y no eval\'ua el
s\'{\i}mbolo como variable de programa.  Su papel es an\'alogo al de un
alias de tabla en una sentencia SQL \texttt{SELECT}.

% ---------------------------------------------------------------------------
\section{Formato condicional declarativo}
\label{sec:formato-condicional}
% ---------------------------------------------------------------------------

El coloreado de celdas es una herramienta habitual en las hojas de c\'alculo
de planificaci\'on: resaltar en rojo los profesores con conflicto de horario
permite detectar problemas de un vistazo.  El DSL ofrece la forma
\texttt{conditional-rendering} para declarar ese coloreado sin escribir
f\'ormulas Excel ni llamadas a la API de formato.

\begin{lstlisting}[language=CommonLisp,
  caption={Sintaxis de conditional-rendering.}]
(conditional-rendering
  :condition      (expresion-booleana)
  :target-columns (col1 col2 ...))
\end{lstlisting}

Cuando el nodo condici\'on es una instancia de \texttt{xl-expr-exists}, el
backend no produce colores est\'aticos en tiempo de compilaci\'on, sino una
\texttt{FormulaRule} de Excel incrustada en el archivo \texttt{.xlsx}.
Esta regla persiste en el libro y se reevalua autom\'aticamente cada vez
que el usuario edita los datos directamente en Excel, sin necesidad de
volver a ejecutar el DSL.

La decisi\'on de usar \texttt{FormulaRule} en lugar de colorear celdas
directamente es exclusiva del backend: el programador del DSL no indica en
ning\'un lugar qu\'e mecanismo de formato se utilizar\'a.

% ............................................................................
\subsection{F\'ormula generada para la variante same-table}
\label{sec:cf-same-table}
% ............................................................................

Cuando \texttt{exists} usa \texttt{:all-rows}, el backend genera una formula
SUMPRODUCT que, para cada fila, multiplica tres factores: coincidencia en
las columnas de agrupaci\'on (\texttt{:matching}), exclusi\'on de la fila
actual (para no comparar una fila consigo misma) y solapamiento en al menos
una columna de rol (\texttt{:any-overlap}).

Para una regla con \texttt{:matching (dia hora)} y
\texttt{:any-overlap (tutor \ldots secretario)}, el backend genera la
siguiente formula, aplicada al rango de la tabla:

\begin{lstlisting}[style=pythonstyle,
  caption={Formula Excel CF para la variante same-table (escrita para la fila 4).}]
=SUMPRODUCT(
    ($G$4:$G$6=$G4) *            -- dia coincide
    ($H$4:$H$6=$H4) *            -- hora coincide
    (ROW($B$4:$B$6)<>ROW($B4)) * -- excluir la fila actual
    (($B$4:$B$6=$B4) + ($C$4:$C$6=$C4) + ... + ($F$4:$F$6=$F4))
) > 0
\end{lstlisting}

Las referencias con columna fija y fila relativa (p.\,ej.\ \texttt{\$G4})
hacen que la f\'ormula se eval\'ue correctamente para cada fila del rango:
al desplazarse de la fila~4 a la fila~5, las referencias de ``fila actual''
avanzan autom\'aticamente con la celda.

% ............................................................................
\subsection{F\'ormula generada para la variante cross-table}
\label{sec:cf-cross-table}
% ............................................................................

Cuando \texttt{exists} usa \texttt{:rows-of}, el backend genera una formula
SUMPRODUCT que combina dos condiciones: que el nombre del profesor aparezca
en al menos una columna de rol de la tabla de defensas (\texttt{:self-in}),
y que ese slot de horario tenga m\'as de una defensa programada
(\texttt{:matching} con \texttt{COUNTIFS}).

Para la regla de \texttt{profesores-table} con
\texttt{:self-in (tutor \ldots secretario)} y \texttt{:matching (dia hora)},
el backend genera la siguiente formula, aplicada al rango de la columna
\texttt{nombre}:

\begin{lstlisting}[style=pythonstyle,
  caption={Formula Excel CF para la variante cross-table (escrita para K4).}]
=SUMPRODUCT(
    (($B$4:$B$6=$K4) + ($C$4:$C$6=$K4) + ($D$4:$D$6=$K4) +
     ($E$4:$E$6=$K4) + ($F$4:$F$6=$K4)) *   -- aparece en algun rol
    (COUNTIFS($G$4:$G$6,$G$4:$G$6,
              $H$4:$H$6,$H$4:$H$6) > 1)      -- slot con mas de 1 defensa
) > 0
\end{lstlisting}

La referencia \texttt{\$K4} corresponde al nombre del profesor en la celda
actual (columna K bloqueada, fila relativa).  Los rangos
\texttt{\$B\$4:\$B\$6} a \texttt{\$F\$4:\$F\$6} son las columnas de rol
de \texttt{defensas-table}, completamente absolutas.  \texttt{COUNTIFS}
cuenta cu\'antas filas comparten el mismo \texttt{dia} y \texttt{hora}; si
el resultado es mayor que~1, ese slot tiene conflicto.

El resultado es que los profesores marcados en rojo son exactamente aquellos
que aparecen en al menos una defensa cuyo slot horario tiene m\'as de una
defensa programada.

% ---------------------------------------------------------------------------
\section{Ensamblaje del workbook}
\label{sec:workbook}
% ---------------------------------------------------------------------------

Una vez definidas las plantillas de tabla, el programador ensambla el
workbook con tres macros: \texttt{libro} crea el nodo ra\'{\i}z del AST con
el nombre del archivo de salida; \texttt{hoja} agrupa tablas y expresiones
fijas en una pesta\~na del libro; \texttt{tabla} instancia una plantilla
(\texttt{def-table}) con datos concretos.

\begin{lstlisting}[language=CommonLisp,
  caption={Ensamblaje del workbook con libro, hoja y tabla.}]
(libro nombre-libro
  :filename "Archivo.xlsx"
  :hojas (list
    (hoja "NombreHoja"
      (tabla plantilla1 :data *DATOS-1*)
      (tabla plantilla2 :data *DATOS-2*)
      :fixed-expressions
        (((nav ...) (expresion))
         ((nav ...) (expresion))))))
\end{lstlisting}

Cuando una \texttt{hoja} recibe m\'as de una \texttt{tabla}, el generador
las coloca en la misma \emph{regi\'on horizontal}: las tablas se disponen de
izquierda a derecha separadas por una columna en blanco.  Este dise\~no
permite que las referencias cross-table (\texttt{trange}) se resuelvan
dentro de la misma regi\'on con letras de columna absolutas correctas.

La clave \texttt{:fixed-expressions} acepta una lista de pares
\texttt{(posici\'on, expresi\'on)}.  La posici\'on se especifica con
\texttt{nav} (secci\'on~\ref{sec:nav}); la expresi\'on puede ser cualquier
forma del DSL.  Estas expresiones quedan fuera de las tablas y su posici\'on
se recalcula seg\'un el tama\~no real de los datos, de modo que los totales
de pie de tabla se desplazan autom\'aticamente si se a\~naden m\'as filas.

% ---------------------------------------------------------------------------
\section{Posicionamiento relativo: \texttt{nav}}
\label{sec:nav}
% ---------------------------------------------------------------------------

Las expresiones fijas de una hoja (etiquetas de totales, celdas de resumen)
no pueden tener coordenadas absolutas codificadas, pues esas coordenadas
dependen de cu\'antas filas de datos tenga cada tabla en cada ejecuci\'on.
La forma \texttt{nav} resuelve esto: calcula la posici\'on de una celda como
un \emph{ancla} fija en una tabla m\'as una secuencia de pasos
direccionales.

Las anclas disponibles son \texttt{ultima-fila} y \texttt{primera-fila},
que identifican la \'ultima o primera fila de datos de una columna de una
tabla.  A partir del ancla se pueden encadenar pasos \texttt{abajo},
\texttt{arriba}, \texttt{izquierda} y \texttt{derecha} con su cantidad:

\begin{lstlisting}[language=CommonLisp,
  caption={Sintaxis de nav y anclas de tabla.}]
;; Una fila debajo de la ultima fila de datos
(nav (ultima-fila tabla-id col-name) abajo 1)

;; Una fila debajo y dos columnas a la derecha
(nav (ultima-fila tabla-id col-name) abajo 1 derecha 2)

;; Una fila encima de la primera fila de datos
(nav (primera-fila tabla-id col-name) arriba 1)
\end{lstlisting}

El compilador resuelve \texttt{ultima-fila} a la fila Excel correspondiente
en tiempo de generaci\'on y produce una referencia de celda absoluta que
siempre apunta al lugar correcto, independientemente del n\'umero de filas
de datos.

\begin{lstlisting}[language=CommonLisp,
  caption={Expresiones fijas de totales con posicionamiento relativo.}]
;; Etiqueta una fila debajo de la ultima defensa
((nav (ultima-fila defensas-table estudiante) abajo 1)
 (str "Total defensas:"))
;; Valor COUNTA en la columna siguiente
((nav (ultima-fila defensas-table estudiante) abajo 1 derecha 1)
 (counta (trange defensas-table estudiante)))

;; Suma de defensas por profesor
((nav (ultima-fila profesores-table nombre) abajo 1 derecha 2)
 (sum-range (trange profesores-table defensas)))
\end{lstlisting}

% ---------------------------------------------------------------------------
\section{Generaci\'on y ejecuci\'on}
\label{sec:generacion}
% ---------------------------------------------------------------------------

Las dos \'ultimas formas del DSL cierran el ciclo de compilaci\'on:
\texttt{xl-generate} recorre el \'arbol AST y produce el programa de
construcci\'on del libro; \texttt{xl-run-generated} ejecuta ese programa
y crea el archivo \texttt{.xlsx}.

\begin{lstlisting}[language=CommonLisp,
  caption={Compilacion del AST y ejecucion.}]
;; Recorre el AST y genera el programa de construccion del libro
(xl-generate horario-tesis "horario-tesis.py")

;; Ejecuta el programa generado y crea el archivo .xlsx
(xl-run-generated "horario-tesis.py")
\end{lstlisting}

El archivo \texttt{.xlsx} resultante contiene tanto los valores y f\'ormulas
de celda como las \texttt{FormulaRule} de formato condicional.  Estas
\'ultimas permanecen activas cuando el usuario abre y edita el libro
directamente en Excel: el coloreado se recalcula sin necesidad de volver a
ejecutar el DSL.

% ---------------------------------------------------------------------------
\section{Ejemplo completo: Horario de Defensas de Tesis}
\label{sec:ejemplo}
% ---------------------------------------------------------------------------

El caso de uso principal del DSL es la generaci\'on autom\'atica de un
horario de defensas de tesis con dos tablas relacionadas.  La primera tabla,
\texttt{defensas-table}, registra cada defensa con su tribunal (cinco roles:
tutor, oponente, presidente, vocal y secretario), su slot horario (d\'{\i}a
y hora) y el local.  La segunda tabla, \texttt{profesores-table}, enumera
los profesores con su grado acad\'emico y el n\'umero de defensas en que
participan; adem\'as, sus filas se colorean en rojo cuando el profesor tiene
un conflicto de horario, es decir, cuando aparece en dos defensas
programadas en el mismo slot.

El listado siguiente es el programa DSL completo.  En \'el se pueden
observar todas las construcciones descritas en el cap\'{\i}tulo: definici\'on
de plantillas de tabla, columna \texttt{computed}, cuantificador
\texttt{exists} cross-table, formato condicional declarativo, ensamblaje del
workbook y posicionamiento relativo con \texttt{nav}.

\begin{lstlisting}[language=CommonLisp,
  caption={Programa DSL completo: Horario de Defensas de Tesis.},
  label={lst:tutorial}]
(load "dsl-directo.lisp")
(load "data-tutorial.lisp")   ; *DATOS-DEFENSAS*, *DATOS-PROFESORES*

;; -- TABLA 1: defensas ------------------------------------------------
;; Nueve columnas: cinco de rol (tutor...secretario) mas dia, hora, local.
;; Las cinco columnas de rol son contiguas para poder referenciarlas como
;; rango con (trange defensas-table tutor secretario).
;; No tiene regla de coloreado; el analisis se delega a profesores-table.
(def-table defensas-table
  ((estudiante "") (tutor "") (oponente "") (presidente "")
   (vocal "") (secretario "") (dia "") (hora "") (local ""))
  :cell-height 1
  :cell-width 1)

;; -- TABLA 2: profesores ----------------------------------------------
;; La columna "defensas" cuenta apariciones del nombre en los cinco roles.
;; El formato condicional marca en rojo a quienes tienen conflicto.
(def-table profesores-table
  ((nombre "") (grado "") (defensas ""))
  :cell-height 1
  :cell-width  1
  :computed
    ((defensas (countif (trange defensas-table tutor secretario)
                        (col nombre))))
  :render
    (conditional-rendering
      :condition
        (exists (r :rows-of defensas-table)
          :self-in  (tutor oponente presidente vocal secretario)
          :matching (dia hora))
      :target-columns (nombre)))

;; -- WORKBOOK ---------------------------------------------------------
(libro horario-tesis
  :filename "Horario_Tesis.xlsx"
  :hojas (list
    (hoja "Defensas"
      (tabla defensas-table   :data *DATOS-DEFENSAS*)
      (tabla profesores-table :data *DATOS-PROFESORES*)
      :fixed-expressions
        (((nav (ultima-fila defensas-table estudiante) abajo 1)
          (str "Total defensas:"))
         ((nav (ultima-fila defensas-table estudiante) abajo 1 derecha 1)
          (counta (trange defensas-table estudiante)))
         ((nav (ultima-fila profesores-table nombre) abajo 1)
          (str "Total profesores:"))
         ((nav (ultima-fila profesores-table nombre) abajo 1 derecha 2)
          (sum-range (trange profesores-table defensas)))))))

(xl-generate horario-tesis "horario-tesis.py")
(xl-run-generated "horario-tesis.py")
\end{lstlisting}

Los datos de entrada contienen tres defensas y seis profesores.  Dos de las
defensas comparten el slot \texttt{Lunes 10:00} con roles solapados.
El programa genera \texttt{Horario\_Tesis.xlsx} con la estructura siguiente:

\begin{itemize}
  \item Columnas A--I: \texttt{defensas-table} con tres filas de datos.
  \item Columna J: separador en blanco de la regi\'on.
  \item Columnas K--M: \texttt{profesores-table} con seis filas de datos
    y la columna \texttt{Defensas} calculada mediante \texttt{COUNTIF}.
  \item Fila de totales debajo de cada tabla: conteo de defensas con
    \texttt{COUNTA} y suma de defensas por profesor con \texttt{SUM}.
  \item Formato condicional \texttt{FormulaRule} aplicado al rango de la
    columna \texttt{nombre}: los profesores que participan en al menos una
    defensa con conflicto de horario aparecen resaltados en rojo.
\end{itemize}

La regla de coloreado generada implementa exactamente la sem\'antica del
\texttt{exists} declarado en el DSL:

\begin{lstlisting}[style=pythonstyle,
  caption={FormulaRule SUMPRODUCT+COUNTIFS incrustada en el .xlsx.}]
-- Rango: K4:K9   Formula escrita para K4 (fila relativa)
=SUMPRODUCT(
    (($B$4:$B$6=$K4) + ($C$4:$C$6=$K4) + ($D$4:$D$6=$K4) +
     ($E$4:$E$6=$K4) + ($F$4:$F$6=$K4)) *
    (COUNTIFS($G$4:$G$6,$G$4:$G$6,
              $H$4:$H$6,$H$4:$H$6) > 1)
) > 0
\end{lstlisting}

La primera parte comprueba si el nombre del profesor (celda \texttt{K4})
aparece en alguna de las cinco columnas de rol de \texttt{defensas-table};
la segunda parte comprueba si ese slot tiene m\'as de una defensa
programada.  Al ser una \texttt{FormulaRule} incrustada en el archivo Excel,
cualquier edici\'on posterior de los datos provoca una reevaluaci\'on
autom\'atica del coloreado, sin necesidad de volver a ejecutar el DSL.

\end{document}
